\documentclass[11pt]{article}
\usepackage{amsmath}
\usepackage[hidelinks]{hyperref}
\title{Staggered notices can make feasible formation endpoints unreachable}
\date{27 November 2026}
\begin{document}
\maketitle
\paragraph{Audit.}
Five homogeneous robots occupy two contiguous two-gap blocks.
All start and finish at rest, \(a=2\), speed cap 2, deadline
\(T=2\). The initial gap is \(g>0\) and every actual gap must
remain in \([g,g+3]\) throughout. One of the two blocks needs
extra terminal span \(\delta\). The middle robot learns which
block at time 0. Its two immediate neighbors learn at time 1;
the two outer robots never do. Until its own receipt each
robot has identical admissible observations in the two worlds,
including no indirect reading of the middle's moving response.
These are given receipt cuts, not a measured network or sensor.
There is room for the described motions at both ends.

\paragraph{The terminal screen accepts four.}
Use displacements, in order from left to right:
\[
       (-2,-1,2,2,2)\quad\hbox{or}\quad(-2,-2,-2,1,2)
\]
for left and right demand respectively. These vectors meet the
terminal gap bounds, both outer robots have common targets, and
the needed block gains 4. They also meet the strongest individual
travel and branch-separation bounds: each absolute displacement
is at most \(B(2)=2\); an immediate neighbor's half-difference is
at most \(B(1)=1/2\). Checking these facts plus all hard terminal
constraints as a difference system would accept the target.
That system is a relaxation here, not a matched sufficient
trajectory test.

\paragraph{A hard gap must be feasible before the deadline too.}
Write \(M_{\rm L}(t),M_{\rm R}(t)\) for the middle robot's
displacements in the two worlds, and \(N_{\rm L},N_{\rm R}\)
for those of its left neighbor. Put
\(D_M=(M_{\rm L}-M_{\rm R})/2\), \(A_M=(M_{\rm L}+M_{\rm R})/2\),
and similarly \(D_N,A_N\).
Let \(F\) integrate the full rest-to-rest velocity envelope
over \([0,2]\), and \(G\) integrate the shorter envelope over
\([1,2]\), extended by zero beforehand. On \([3/2,2]\),
\(F-G=3/2\). If both demands are met with gain
\(\delta=4-\eta\), necessarily \(\eta\ge0\) and
\(D_M(2)\ge 2-\eta\), by adding the two terminal demand
inequalities and bounding the two outside displacements.
At every point the middle's half-difference velocity is
bounded above by the full velocity envelope \(u\);
also \(|A_M'|\le u-D_M'\), by the two separate bounds
on \(|M_{\rm L}'|\) and \(|M_{\rm R}'|\).
The nonnegative deficit integrated to time 2 is at most
\(\eta\). Thus \(D_M(t)\ge F(t)-\eta\) and
\(|A_M(t)|\le\eta\) for every \(t\).
The neighbor's half-difference starts at rest at its
time-1 notice, so \(D_N(t)\le G(t)\).

On the entire final half-unit of time, the half-difference in
the left near gap is consequently at least \(3/2-\eta\).
Both real gaps must stay in the same allowed interval of
width 3. Their mean extra gap \(A_M-A_N\) lies
in \([3/2-\eta,3/2+\eta]\). In particular
\[
     A_N(t)\in[-3/2-2\eta,\,-3/2+2\eta]
       \qquad(3/2\le t\le2).
\]
At exact terminal saturation \(\eta=0\), this forces the
left neighbor's mean to move by \(3/2\) and become stationary
by time \(3/2\). The mean is itself acceleration bounded,
and its maximum rest-to-rest travel in that time is \(9/8\).
So the terminal target 4 is physically unreachable.

The obstruction persists below 4. Let
\(q=|A_N'(3/2)|\). Since the mean is at rest at time 2,
maximal braking takes at most the remaining half-unit.
The unavoidable displacement before stopping from speed \(q\)
is at least \(q^2/4\), but the final band's full width is
\(4\eta\). Hence \(q^2\le16\eta\).
From rest to time \(3/2\) with final speed magnitude \(q\),
the pointwise envelope
\(\min(2t,q+2(3/2-t))\) bounds the absolute velocity.
Its integral is at most \(9/8+3q/4\).
Since at least \(3/2-2\eta\) was required, it follows that
\[
             3/8\le 2\eta+3q/4\le2\eta+3\sqrt\eta .
\]
Consequently \(\eta\ge ((2\sqrt3-3)/4)^2\). This is an
explicit strict upper bound, not a claim of tightness.

\paragraph{Positive matched trajectory.}
In the left-demand world take the five displacements
\[
 \left(-F,\;-\tfrac34F+\tfrac14G,\;
             \tfrac34F+\tfrac14G,\;
             \tfrac34F+\tfrac14G,\;F\right)
\]
and mirror them in the right-demand world.
Each robot uses profiles with coefficient magnitudes
summing to at most 1. Therefore the original acceleration
and speed bounds hold. The neighbors' sign-dependent term
is zero until time 1, and the outer motions are common.
The left-world extra gaps are
\((F+G)/4,\ 3F/2,\ 0,\ (F-G)/4\),
all in \([0,3]\) for the whole motion. The target block
gains \(29/8\), greater than the middle-only capacity
\(7/2\) at the given cut. This is a constructive lower
bound and an explicit example where a late neighbor
receipt helps. It does not reach the relaxed terminal
upper value.

\paragraph{Verification and residual.}
The proof separates a false positive in a terminal-only
relaxation from an independently feasible continuous
policy. Four Lean lemmas check the real algebra, and
4,097 exact dyadic instants check the attainable profile,
its evidence cuts, and the saturated envelope comparison.
There is no assertion that the grid is a runtime safety
test or that the upper and lower bounds identify the exact
optimum. The result does not re-open the equal-cut theorem
of PR~712, and no actual observation, independent mission,
or R6 acceptance is claimed.
\end{document}
